\chapter{Algorithm Comparison}
\label{chap:comparison}

Both IRLS (Chapter~\ref{chap:irls}) and the Deflected Subgradient Method
(Chapter~\ref{chap:dsm}) solve the same LASSO problem
\begin{equation}
    \min_{\mathbf{w} \in \mathbb{R}^H} f(\mathbf{w})
    = \|\mathbf{X}\mathbf{w} - \mathbf{y}\|_2^2 + \lambda\|\mathbf{w}\|_1,
    \label{eq:comp_objective}
\end{equation}
yet they differ fundamentally in \emph{how} they handle the non-smoothness of the
$\ell_1$ penalty and in the trade-offs that result.
This chapter compares the two methods along the dimensions that matter most in practice:
mathematical strategy, convergence behaviour, computational cost, solution quality,
and suitability for the ELM setting.

% ------------------------------------------------------------------
\section{Core strategy: how each algorithm handles non-smoothness}
\label{sec:comp_strategy}

The root difficulty in~\eqref{eq:comp_objective} is the $\ell_1$ term: it is convex but
non-differentiable wherever any $w_i = 0$, so classical gradient methods cannot be applied
directly.
The two algorithms adopt opposite philosophies.

\paragraph{IRLS — smoothing via majorization.}
IRLS sidesteps the non-smoothness by never minimising $f$ directly.
At each iteration, it constructs a smooth \emph{surrogate} $Q(\mathbf{w}, \mathbf{w}_k)$
that majorizes $f$ (Section~\ref{sec:irls_mm}):
\[
    Q(\mathbf{w}, \mathbf{w}_k)
    = \|\mathbf{X}\mathbf{w} - \mathbf{y}\|_2^2
    + 2\lambda\mathbf{w}^T\mathbf{W}_k\mathbf{w},
\]
where the non-smooth $\|\mathbf{w}\|_1$ has been replaced by the quadratic
$\mathbf{w}^T\mathbf{W}_k\mathbf{w}$, with weights
$(\mathbf{W}_k)_{ii} = \max(|w_{k,i}|,\varepsilon_{\mathrm{thr}})^{-1/2}$ that encode
the current estimate of the solution's sparsity pattern (Section~\ref{sec:irls_idea}).
Minimising $Q(\cdot, \mathbf{w}_k)$ amounts to solving a weighted $L_2$-regularised
least-squares problem with a closed-form solution via the normal equations
(Section~\ref{sec:irls_subproblem}).

\paragraph{DSM — direct subgradient with memory.}
DSM confronts the non-smoothness head-on. It descends along a \emph{subgradient}
$\tilde{\mathbf{g}} \in \partial f(\mathbf{w}_i)$ at every step, computing it exactly
via the sum rule $\tilde{\mathbf{g}} = 2\mathbf{X}^T(\mathbf{X}\mathbf{w}_i-\mathbf{y})
+ \lambda\mathbf{s}$ (Section~\ref{sec:dsm_subgrad}).
To combat the zig-zag oscillations of the plain subgradient method
(Section~\ref{sec:dsm_motivation}), the direction is \emph{deflected}: it blends the
current subgradient with the previous direction (Section~\ref{sec:dsm_deflection}), and
the Polyak stepsize is controlled via a target-level estimate of $f^*$
(Section~\ref{sec:dsm_target_level}).

\paragraph{Key structural difference.}
IRLS transforms a non-smooth problem into a sequence of smooth ones, then solves each
exactly (via a linear system solve). DSM never smooths the objective; instead it relies
on the convergence theory for subgradient methods, accepting that individual steps may
increase $f$.

% ------------------------------------------------------------------
\section{Convergence behaviour}
\label{sec:comp_convergence}

The convergence properties of the two algorithms differ in both \emph{rate} and
\emph{monotonicity}.

\subsection{Monotonicity}

IRLS is \textbf{monotone}: the MM framework guarantees
$f(\mathbf{w}_{k+1}) \leq f(\mathbf{w}_k)$ at every iteration
(Section~\ref{sec:irls_mm}).
This makes convergence easy to monitor and diagnose — a non-decreasing plot of
$f(\mathbf{w}_k)$ immediately signals a bug.

DSM is \textbf{non-monotone}: individual function values $f(\mathbf{w}_i)$ oscillate
throughout the run (Section~\ref{sec:dsm_convergence}).
The meaningful quantity is the \emph{record value}
$\bar{f}^{\,i} = \min_{h \leq i} f(\mathbf{w}_h)$, which decreases by construction
but may improve only sporadically.
Plotting $f(\mathbf{w}_i)$ without tracking $\bar{f}^{\,i}$ gives a misleading picture.

\subsection{Rate}

IRLS converges at a \textbf{linear} (geometric) rate: on a semilogarithmic plot of
$f(\mathbf{w}_k) - f^*$ versus iteration number, the curve is approximately a straight
line (Section~\ref{sec:irls_convergence}).
In practice, 10--50 iterations suffice to reach high accuracy.

DSM converges at a \textbf{sublinear} rate, $O(1/\varepsilon^2)$
(Theorem~\ref{thm:dsm_convergence}): on the same semilogarithmic plot, the curve
flattens progressively.
This rate is information-theoretically optimal for nonsmooth convex functions
(Section~\ref{sec:dsm_convergence}), so the gap is intrinsic to the approach, not an
implementation defect. Concretely, achieving $\varepsilon = 10^{-4}$ requires roughly
$100\times$ more iterations than $\varepsilon = 10^{-3}$.

\subsection{Convergence on the ELM problem}

For the ELM LASSO problem, $\mathbf{X} = \sigma(\mathbf{X}_{\text{origin}}\mathbf{W}_1)$
has full column rank with high probability (since $\mathbf{W}_1$ is random and $\sigma$ is
a nonlinear activation), so $f$ is strictly convex and the unique minimiser $\mathbf{w}^*$
exists.
This means IRLS converges to $\mathbf{w}^*$ (Theorem~\ref{thm:irls_convergence}) with
a linear rate governed by the condition number of the weighted normal-equation matrices.
DSM also converges to $\mathbf{w}^*$ (Theorem~\ref{thm:dsm_convergence}), and the
strong convexity of $f$ improves the theoretical lower bound from $O(1/\varepsilon^2)$
to $O(1/\varepsilon)$, though the algorithm does not exploit this structure.

% ------------------------------------------------------------------
\section{Computational cost}
\label{sec:comp_cost}

\subsection{Per-iteration cost}

\begin{center}
\begin{tabular}{lll}
\toprule
\textbf{Operation} & \textbf{IRLS} & \textbf{DSM} \\
\midrule
Weight / direction update & $O(H)$ & $O(H)$ \\
Matrix-vector product(s) & --- & $O(MH)$ \\
Linear system solve & $O(H^3)$ (Cholesky) or $O(kH^2)$ (CG) & --- \\
\midrule
\textbf{Total per iteration} & $O(H^3)$ & $O(MH)$ \\
\bottomrule
\end{tabular}
\end{center}

IRLS requires solving an $H \times H$ symmetric positive definite linear system at every
iteration (Section~\ref{sec:irls_subproblem}).
With Cholesky factorisation the cost is $O(H^3/3)$ per iteration.
With $k$ steps of Conjugate Gradient it is $O(kH^2)$, preferable when $H$ is large and
$\mathbf{Q}_k$ is well-conditioned.
The matrices $\mathbf{A} = \mathbf{X}^T\mathbf{X}$ and $\mathbf{b} = \mathbf{X}^T\mathbf{y}$
are pre-computed once at cost $O(MH^2)$ and $O(MH)$ respectively.

DSM requires only two matrix-vector products per iteration (to compute the subgradient:
$\mathbf{X}\mathbf{w}_i$ costs $O(MH)$ and $\mathbf{X}^T(\cdot)$ costs $O(MH)$), giving
a total per-iteration cost of $O(MH)$ (Section~\ref{sec:dsm_subgrad}).
No precomputation of $\mathbf{X}^T\mathbf{X}$ is needed.

\subsection{Total cost}

Let $k_{\mathrm{IRLS}}$ and $k_{\mathrm{DSM}}$ denote the respective iteration counts.
\begin{align*}
    \text{IRLS total:} &\quad O(MH^2) + k_{\mathrm{IRLS}} \cdot O(H^3), \\
    \text{DSM total:}  &\quad k_{\mathrm{DSM}} \cdot O(MH).
\end{align*}
Since IRLS converges linearly and DSM converges sublinearly, the iteration counts satisfy
$k_{\mathrm{IRLS}} \ll k_{\mathrm{DSM}}$ for a given accuracy.
The regime where DSM is competitive is therefore when $H$ is large enough that
$H^3 \gg MH$, i.e.\ $H^2 \gg M$, so that the Cholesky cost of IRLS dominates.

For the ELM problem, $M$ is the number of training samples and $H$ is the hidden layer
size. Typical configurations have $M \gg H$ (many samples, moderate hidden layer), so
$O(MH^2) \gg O(H^3)$ and the pre-computation of $\mathbf{A}$ dominates IRLS's total
cost.
In this regime, the total costs are roughly $O(MH^2)$ for IRLS and
$O(k_{\mathrm{DSM}} \cdot MH)$ for DSM — IRLS wins whenever
$H \lesssim k_{\mathrm{DSM}}$, which holds for any reasonable accuracy requirement.

% ------------------------------------------------------------------
\section{Solution quality}
\label{sec:comp_quality}

\subsection{Sparsity}

A central goal of LASSO is to produce a \emph{sparse} solution
$\mathbf{w}^*$ with many components exactly zero, which in the ELM context corresponds
to identifying the most relevant hidden neurons and discarding the rest.

IRLS promotes \textbf{exact sparsity}: as iterations proceed, components $w_{k,i}$ with
small magnitude receive increasingly large weights $(\mathbf{W}_k)_{ii}$, pushing them
further toward zero.
At convergence, many components are exactly zero (within $\varepsilon_{\mathrm{thr}}$),
yielding a genuinely sparse weight vector (Section~\ref{sec:irls_expected}).
The sparsity structure emerges naturally from the algorithm's geometry.

DSM produces \textbf{approximate sparsity}: the iterates drift near zero for irrelevant
components, but they are rarely driven to exact zero because the subgradient steps carry
no mechanism that enforces zeros.
A post-processing thresholding step (setting $w_i \gets 0$ if $|w_i| < \tau$ for some
small $\tau$) is typically needed to obtain a truly sparse solution
(Section~\ref{sec:dsm_properties}).

\subsection{Accuracy}

For a fixed computational budget (number of iterations or wall-clock time), IRLS
reaches substantially higher accuracy than DSM on small-to-moderate problems, due to
its linear convergence rate. DSM can reach competitive accuracy when the per-iteration
cost advantage of $O(MH)$ versus $O(H^3)$ outweighs the slower convergence rate.

% ------------------------------------------------------------------
\section{Parameter sensitivity}
\label{sec:comp_params}

\paragraph{IRLS.}
The algorithm has two main hyperparameters: the regularisation strength $\lambda$ (shared
with the problem formulation) and the safety threshold $\varepsilon_{\mathrm{thr}}$.
The latter controls the trade-off between numerical stability and approximation quality:
too large degrades sparsity, too small risks division by zero.
A value around $10^{-8}$ is a robust default.
The convergence stopping criterion $\varepsilon_{\mathrm{stop}}$ is straightforward to
tune: the relative change $\|\mathbf{w}_{k+1}-\mathbf{w}_k\|/\max(1,\|\mathbf{w}_k\|)$
provides a scale-invariant and interpretable criterion (Section~\ref{sec:irls_algorithm}).

\paragraph{DSM.}
The algorithm involves four additional hyperparameters: $\delta_0$, $\rho$, $R$, and
$\beta$ (Section~\ref{sec:dsm_parameters}).
Their calibration is critical and interdependent: a poor choice of $\delta_0$ dominoes
into poor target levels, which cause either oversized steps ($\delta_0$ too large) or
stalling ($\delta_0$ too small).
The patience threshold $R$ and reduction rate $\rho$ control adaptivity.
Unlike IRLS, there is no reliable stopping criterion based on iterate change, since the
subgradient norm may remain large even at the optimum (Section~\ref{sec:dsm_properties}).
A fixed iteration budget is the most practical approach.

% ------------------------------------------------------------------
\section{Summary and guidelines for the ELM problem}
\label{sec:comp_summary}

\begin{center}
\begin{tabular}{lll}
\toprule
\textbf{Criterion} & \textbf{IRLS} & \textbf{DSM} \\
\midrule
Approach & Smoothing (MM) & Direct subgradient \\
Monotone & Yes & No (record value only) \\
Convergence rate & Linear & Sublinear $O(1/\varepsilon^2)$ \\
Per-iteration cost & $O(H^3)$ & $O(MH)$ \\
Precomputation & $O(MH^2)$ & None \\
Exact sparsity & Yes & No (needs thresholding) \\
Parameters to tune & $\varepsilon_{\mathrm{thr}}$ & $\delta_0, \rho, R, \beta$ \\
Stopping criterion & Reliable (iterate change) & Unreliable (fixed budget) \\
Preferred when & High accuracy, $H \ll M$ & Large $H$, moderate accuracy \\
\bottomrule
\end{tabular}
\end{center}

\paragraph{When to prefer IRLS.}
IRLS is the method of choice when:
\begin{itemize}[nosep]
    \item High solution accuracy is required (e.g.\ for precise feature selection or
    reproducibility of sparsity patterns).
    \item $H$ is small to moderate relative to $M$, so the $O(H^3)$ per-iteration
    cost is affordable.
    \item Exact zeros in $\mathbf{w}^*$ are important (e.g.\ to identify which hidden
    neurons to prune).
    \item Easy debugging is preferred: the monotone decrease of $f$ is a strong
    sanity check.
\end{itemize}

\paragraph{When to prefer DSM.}
DSM is competitive when:
\begin{itemize}[nosep]
    \item $H$ is large enough that Cholesky factorisation is infeasible or too slow
    ($H^2 \gg M$).
    \item Moderate accuracy ($\varepsilon \approx 10^{-3}$) suffices.
    \item Memory is a constraint: DSM does not require storing $\mathbf{X}^T\mathbf{X}
    \in \mathbb{R}^{H \times H}$, only the current iterate and direction.
    \item The problem changes online (e.g.\ streaming data), where recomputing
    $\mathbf{A}$ is expensive.
\end{itemize}

\paragraph{Practical verdict for the ELM problem.}
In the typical ELM setting (fixed random $\mathbf{W}_1$, offline training), $M$ is large
and $H$ is chosen much smaller, so $M \gg H$.
In this regime IRLS is clearly preferable: the precomputed $\mathbf{A} =
\mathbf{X}^T\mathbf{X}$ is computed once at cost $O(MH^2)$, and each subsequent
iteration costs only $O(H^3)$, which is modest for typical $H \leq 1000$.
The monotone convergence, reliable stopping criterion, exact sparsity, and minimal
parameter tuning all favour IRLS.
DSM serves as a useful baseline and as a fallback for very large $H$, where it
trades accuracy for scalability.
